\begin{savequote}[8cm]
\textlatin{Neque porro quisquam est qui dolorem ipsum quia dolor sit amet, consectetur, adipisci velit...}

There is no one who loves pain itself, who seeks after it and wants to have it, simply because it is pain...
  \qauthor{--- Cicero's \textit{de Finibus Bonorum et Malorum}}
\end{savequote}

\chapter{\label{ch:1-intro}Introduction} 

\minitoc

\section{Motivation}

Quantum computing, a field rooted in the principles of quantum mechanics, represents a revolutionary approach to computation. 
Unlike classical computing, which relies on binary operations performed on bits, 
quantum computing utilizes quantum bits, or qubits, capable of representing and processing information in ways that classical bits cannot. 
This unique capability of qubits to exist in superpositions and entangle with each other allows quantum computers 
to perform certain types of calculations exponentially faster than their classical counterparts.
If a classical computer were limited by memory and runtime, 
a quantum computer could potentially bypass these scaling issues - dubbed the "quantum advantage".

If scalable quantum computers were to become a reality, 
they could solve some of the most challenging problems in computer science much more efficiently than classical computers. 
One notable example is Shor's algorithm for factoring large numbers, which could break widely-used cryptographic systems such as RSA. 
Another significant algorithm is Grover's search algorithm, which provides quadratic speed-ups for unstructured search problems. 

As we strive to advance the field of quantum computing, 
classical computations play an essential role in assisting and accelerating the development of quantum computers. 
Classical simulation is a fundamental aspect of understanding and designing quantum hardware, 
often serving as the only means to validate these quantum systems \cite{xu2023herculean}. 
Additionally, quantum algorithms can be implemented on classical simulators for testing and verification purposes 
while full-fledged quantum computers remain beyond reach (for now). 
These simulators emulate the behavior of quantum computers, allowing researchers to simulate processes as if running on actual quantum devices.

% \section{Contribution}

This theses aims to present the speedup of classical simulation of quantum circuits in an almost entirely graphical approach.
From the introduction of the qubits, to the representation of quantum circuits and linear maps,
the underlying linear algebra that governs the operations of quantum computing is represented in a graphical manner.
Using ZX-Calculus, a graphical calculus detailed in \autoref{sec:zx-calculus},
quantum circuits can be represented as diagrams, 
allowing for different discoveries to be more easily made when analysing the structure of said diagrams.

\section{Basic Notation}
\label{sec:basic-notation}
In classical computing, a bit is the fundamental unit of information, 
and is denoted by a boolean value $b\in\{0,1\}$.
In quantum computing, the fundamental unit of information is the qubit,
which we denote as $\ket{b}$, where $b\in\{0,1\}$.
Diagramatically, they are represented as follows:
\[\point{0} \quad\quad\quad\quad \point{1}\]

We consider these the computational or Z-basis states.

\begin{definition}[Z-basis]
  The Z-basis or computational basis is defined as
  \[\left\{\cbox{\point{0}}, \cbox{\point{1}}\right\}\]
\end{definition}

So far, the expressive power is still similar to classical computing.
However, the power of quantum computing comes from the ability to exist in superpositions of these states.
Concretely, the qubit can be represented in a linear combination of the basis states.
\begin{definition}[Pure state]
  A pure state $\ket{\psi}$ is a state
  \[\cbox{\point{\psi}} = a \cbox{\point{0}} + b \cbox{\point{1}}\]
  where $a,b\in\C$ and $|a|^2 + |b|^2 = 1$.
\end{definition}

Then, the '+' or X-basis states can be defined.
\[\cbox{\graypointmap{0}} = \frac{1}{\sqrt{2}}\left(\cbox{\point{0}} + \cbox{\point{1}}\right)\]
\[\cbox{\graypointmap{1}} = \frac{1}{\sqrt{2}}\left(\cbox{\point{0}} - \cbox{\point{1}}\right)\]
Along with the Z-basis states, these states form the building blocks of understanding quantum computing and the use of ZX-calculus for diagrammatic reasoning.

Intuitively, using the diagrammatic representation, we can consider combining diagrams in parallel.
In the literature, this is the same as taking a \textit{tensor product}.
For example, we can have a 2-qubit state.
\[\cbox{\twopoint{10}} = \tikzfig{01state}\]

States can be transformed by quantum gates simply by \textit{composing} them to obtain a new state.
Given that composing the quantum gate $U$ to a state $\ket{\psi}$ gives a state $\ket{\psi'}$,
we can just connect the wires.
\[\cbox{\point{\psi'}} = \cbox{\boxpointmap{\psi}{U}}\]

Intuitively, this means that the wire by itself should just be the identity gate, 
since composing it with any state should not change that state.
\[\cbox{\boxmap{I}} := \cbox{\idwire}\]
\[\cbox{\boxpointmap{\psi}{I}} = \cbox{\longpoint{\psi}} = \cbox{\point{\psi}}\]


The dual of states are the effects. Here, we have the Z-basis effects.
\[\copointmap{0} \quad\quad\quad\quad \copointmap{1}\]

For a state $\ket{\psi}$ and an effect $\bra{\phi}$, we can obtain a scalar inner product.
This is essentially the \textit{bra-ket} notation in diagrammatic form by just connecting the wires, 
which comes with a few other definitions.

\begin{definition}[Inner product \cite{coecke2017picturing}]
  The \textit{inner product} $\braket{\phi}{\psi}$ of two states $\ket{\psi}$ and $\ket{\phi}$ is defined as
  \[\pointbraket{\psi}{\phi}\]
  They are \textit{orthogonal} if 
  \[\cbox{\pointbraket{\psi}{\phi}} = 0\]
  The \textit{squared-norm} of a state $\ket{\psi}$ is
  \[\pointbraket{\psi}{\psi}\]
  $\ket{\psi}$ is \textit{normalised} if 
  \[\cbox{\pointbraket{\psi}{\psi}} = \cbox{\emptydiag} = 1\]
  where the empty diagram $\emptydiag$ represents the scalar 1.
\end{definition}

We can then define the inner product of the Z-basis states.
\begin{definition}
  The inner product of the Z-basis states is
  \[\cbox{\pointbraket{0}{0}} = \cbox{\pointbraket{1}{1}} = \cbox{\emptydiag}\]
  \[\cbox{\pointbraket{1}{0}} = \cbox{\pointbraket{0}{1}} = 0\]
  Succinctly, for $i,j\in\{0,1\}$,
  \[\cbox{\pointbraket{i}{j}} = \delta_{ij} = \begin{cases}
    \cbox{\emptydiag} & \text{if } i = j\\
    0 & \text{if } i \neq j
  \end{cases}\]
  where $\delta_{ij}$ is the Kronecker delta.
\end{definition}
In fact, the succinct definition can be used to determine if a basis is an \textit{orthonormal basis} (ONB).
The X-basis is also an ONB.

We are almost at the stage where we can obtain probabilities from the diagrams.
The scalar inner product is not really a probability, that is,
a number between 0 and 1, but instead, a number $z\in\C$.

First, we note that putting multiple scalars in a single diagram is just the same as multiplying them together.
\[\begin{matrix}
  \pointbraket{\psi}{\phi} \\
  \pointbraket{\psi'}{\phi'}
\end{matrix} =
\cbox{\pointbraket{\psi}{\phi}} \cdot \cbox{\pointbraket{\psi'}{\phi'}}
\]

We also know that for a complex number and its conjugate, $z,\bar{z}\in\C, |z|^2 = z\bar{z}$
Translating that to the diagrammatic representation, we can obtain the following.
\[\left|\cbox{\kpointbraket{\psi}{\phi}}\right|^2 = 
\begin{matrix}
  \kpointbraket{\psi}{\phi}\\
  \kpointbraketconj{\psi}{\phi}
\end{matrix}\]
Note that assymetry is introduced to denote the conjugation. 
Remember that a complex number $z = a + bi \in \C$ has a conjugate $\bar{z} = a - bi \in \C$,
so this corresponds to a vertical reflection in the diagram above.

Using the \textit{Born rule}, if we make sure that $\ket{\psi}$ and $\ket{\phi}$ are normalised,
we can obtain the following.

\[0 \leq \begin{matrix}
  \kpointbraket{\psi}{\phi}\\
  \kpointbraketconj{\psi}{\phi}
\end{matrix}
\leq 1\]

This can be interpreted as our desired probaility - the probability of measuring $\ket{\psi}$ in the state $\ket{\phi}$.

We now have all we need to introduce the ZX-calculus notation.

\section{ZX-Calculus}
\label{sec:zx-calculus}
So far, we have introduced the representation of quantum states, effects, and inner products in a diagrammatic form.
However, there still isn't much we can do without a way to manipulate them.
In the previous section, we mentioned that we can apply quantum gates to states to obtain new states.
Here, we will introduce the ZX-Calculus and show the highly intuitive representation of some quantum gates
such as the Hadamard and CNOT gates, based on the building blocks from the previous section.
The ZX-Calculus is a graphical calculus for quantum computations, developed by Coecke and Duncan in 2008 \cite{Coecke_2011}.
All operations are expressed as \textit{spiders}, connected by \textit{edges} (or \textit{wires}).

\subsection{Spiders}
\begin{definition}[Spiders]
  \[\tikzfig{spider_def_full}\]
\end{definition}
Here, we have green and red spiders, also known as Z and X spiders,
made of the Z and X basis states respectively.
Wires coming in from the left are \textit{inputs}, and wires going out to the right are \textit{outputs}.
The $\alpha$ inside a spider node is also called the \textit{phase} of the spider.
For example, with $\alpha = 0, e^{i\alpha} = \cbox{\emptydiag}$.
It is convenient to denote spiders with $\alpha = 0$ as having no angle written inside the spider.

Immediately, we can see that the Z and X basis states can be formed from the definition of the spiders,
up to the multiplicative constant $\sqrt{2}$.
Note that $e^{i\pi} = -1$.
\begin{equation}
\begin{aligned}
\cbox{\greenbasisspider} &= \cbox{\point{0}} + \cbox{\point{1}} = \sqrt{2} \cbox{\graypointmap{0}}\\
\cbox{\greenpibasisspider} &= \cbox{\point{0}} - \cbox{\point{1}} = \sqrt{2} \cbox{\graypointmap{1}}\\
\cbox{\redbasisspider} &= \cbox{\graypointmap{0}} + \cbox{\graypointmap{1}} = \sqrt{2} \cbox{\point{0}}\\
\cbox{\redpibasisspider} &= \cbox{\graypointmap{0}} - \cbox{\graypointmap{1}} = \sqrt{2} \cbox{\point{1}}
\end{aligned}
\end{equation}

It turns out that in reasoning with spiders, we can ignore the multiplicative constants,
since much of the manipulation of the diagrams is based on the structure of the diagrams themselves.
The scalars are easy to determine using the definitions we have seen so far.
By using $\approx$ to denote equality up to a non-zero factor, we have
\begin{equation}
\begin{aligned}
\cbox{\greenbasisspider} &\approx \cbox{\graypointmap{0}} \quad\quad
\cbox{\greenpibasisspider} \approx \cbox{\graypointmap{1}}\\
\cbox{\redbasisspider} &\approx \cbox{\point{0}} \quad\quad
\cbox{\redpibasisspider} \approx \cbox{\point{1}}
\end{aligned}
\end{equation}

We can also see the single-qubit gates in the ZX-Calculus.
\[\cbox{\phase{\alpha}} 
= \cbox{\pointcopointmap{0}{0}} + e^{i\alpha} \cbox{\pointcopointmap{1}{1}} = Z_{\alpha}\]
\[\cbox{\redphase{\alpha}} 
= \cbox{\graypointcopointmap{0}{0}} + e^{i\alpha} \cbox{\graypointcopointmap{1}{1}} = X_{\alpha}\]

We can then define many of the common quantum gates in this way.
\begin{equation}
\begin{aligned}
  I &= \idphase{} = \redidphase{} = \cbox{\idwire}\\
  X &= \cbox{\redphase{\pi}} \\
  Z &= \cbox{\phase{\pi}} \\
  H &= \tikzfig{had_common} \\
  S &= \cbox{\phase{\frac{\pi}{2}}}\\
  T &= \cbox{\phase{\frac{\pi}{4}}} \\
  CNOT &= \sqrt{2} \cdot \tikzfig{cnot}
\end{aligned}
\end{equation}


There are 2 gates of special note here. 
The Hadamard gate $H$, denoted by the small yellow box, or as a blue dashed line,
is a single-qubit gate that interchanges $Z$ and $X$ bases.
We will see later in the ZX-rules that this is vital for the colour-change rule,
as well as why it is useful to denote it as the blue dashed line.

The CNOT gate is a two-qubit gate that acts on the target qubit if the control qubit 
is in the state $\ket{1} \approx \redpibasisspider$.
The vertical line is used simply because it does not matter which direction the line is going,
where the following equality of the second and third forms can be verified.
\[\tikzfig{cnot} = \tikzfig{cnot-left} = \tikzfig{cnot-right}\]

Last but not least, we also introduce the SWAP gate, which as the name suggests, swaps two qubits.
\[SWAP = \tikzfig{swap} := \sum_{i,j\in \{0,1\}} \cbox{\twotwoketbra{i}{j}{j}{i}} \]

We can now introduce the ZX-rules in the next section.

\subsection{ZX-Rules}
The ZX-Calculus is based on a set of rules that allow us to manipulate the diagrams.
For $\alpha,\beta\in[0,2\pi)$ and $a\in\{0,1\}$, the rules are as follows.

\begin{figure}[h!]
\ctikzfig{ZX-rules}
\caption{$ZX$-rules presented in \cite{Kissinger_2019}}\label{table:rewrite}
\end{figure}

The names of the rules are spider fusion \SpiderRule, Hadamard rule \HadamardRule,
$\pi$-commutation \PiRule, copy rule \CopyRule, bialgebra \SCRule,
identity rule \IdRule, and Hadamard cancel rule \HHRule.
By duality, the rules also hold if the colours are swapped.

The power of ZX-diagrams comes from the fact that two diagrams are equivalent
as long as the connectivity of the diagram is the same -
that is, as long as order of inputs and outputs of the whole diagram is preserved.
This does not depend on the position of the spiders or the direction of wires between them.
This neat feature of ZX-diagrams is called "Only Connectivity Matters" (OCM).

We can then obtain cups and cups using 2-legged spiders together with \IdRule.

\[\tikzfig{cap} \eqrule{\IdRule} \tikzfig{cap_spider} \quad\quad\quad 
\tikzfig{cup} \eqrule{\IdRule} \tikzfig{cup_spider}\]

Caps and cups obey the \textit{yanking equations}.
\[ \tikzfig{line_yank} \qquad\qquad \tikzfig{line_yank2} \]

Before we move on to reasoning about the diagrams in a graph-theoretic manner,
we first introduce the concept of a \textit{Clifford+T diagram}.

\begin{definition}[Clifford diagram]
  A \textit{Clifford diagram} is a ZX-diagram where the phases are restricted to integer multiples of $\frac{\pi}{2}$.
\end{definition}

\begin{definition}[Clifford+T diagram]
  A \textit{Clifford+T diagram} is a ZX-diagram where the phases are restricted to integer multiples of $\frac{\pi}{4}$.
\end{definition}

It is clear from the definitions that a Clifford diagram is a subset of a Clifford+T diagram.
The reason we introduce these concepts is that in trying to obtain our main goal of the thesis,
it is sufficient to reason about Clifford+T diagrams, as a result of the following theorem.

\begin{theorem}[Approximate Universality \cite{Nielsen_Chuang_2010}]
  Clifford+T diagrams are \textit{approximately universal} for quantum computing - 
  they can approximate any quantum circuit.
\end{theorem}

The proof of the above theorem involves the \textit{Solovay-Kitaev theorem},
combining with the fact that any single-qubit phase gate can be approximated by a sequence of $H$ and $T$ gates,
with details found in \cite{Nielsen_Chuang_2010}.
